\documentclass[a4paper, 12pt]{article} %Formato de plantilla que vamos a utilizar
\usepackage[utf8]{inputenc}
\usepackage[spanish]{babel}
\usepackage[a4paper, total={6.5in, 9.5in}]{geometry}
\usepackage{graphicx} % Insertar imagenes
\usepackage[table,xcdraw]{xcolor} % Deteccion de colores
\usepackage{fancyhdr} % Crear cabeceros
\usepackage[colorlinks=true,linkcolor=blue]{hyperref} % Enlances
\usepackage{listings} % Bloques de código

% Declaracion de colores
\definecolor{colorTitulo}{HTML}{0F52BA}
\definecolor{codegreen}{rgb}{0,0.6,0}
\definecolor{codegray}{rgb}{0.5,0.5,0.5}
\definecolor{codepurple}{rgb}{0.58,0,0.82}
\definecolor{backcolour}{rgb}{0.95,0.95,0.92}

% Declaracion de variables
\newcommand{\logoPunta}{img/LogoPuntaDelVerde.png} % Logo del punta
\newcommand{\logoKeter}{img/logo.png} % Logo de la practica
\newcommand{\nombrePractica}{Keter Vulnerability} % Nombre de la practica
\newcommand{\bigpar}{\par\vspace*{0.75cm}}

% Bloques de código
\lstdefinestyle{mystyle}{
    backgroundcolor=\color{backcolour},   
    commentstyle=\color{codegreen},
    keywordstyle=\color{magenta},
    numberstyle=\tiny\color{codegray},
    stringstyle=\color{codepurple},
    basicstyle=\ttfamily\footnotesize,
    breakatwhitespace=false,         
    breaklines=true,                 
    captionpos=b,                    
    keepspaces=true,                 
    numbers=left,                    
    numbersep=5pt,                  
    showspaces=false,                
    showstringspaces=false,
    showtabs=false,                  
    tabsize=2
}

\lstset{style=mystyle}

% Header y footer
\pagestyle{fancy}
\fancyhf{}
\rhead{\includegraphics[width=15cm]{\logoPunta}}
\rfoot{Keter Vulnerability\par Página \thepage}
\lfoot{Emilio Sanchez Garcia\par 2º ASIR}
\renewcommand{\headrulewidth}{0pt}
\renewcommand{\footrulewidth}{1pt}
\setlength\headheight{40.62811pt} % quitar aviso cabecero
\addtolength{\topmargin}{-0.62811pt}

% Comienzo del documento
\begin{document}
% Portada
\begin{titlepage}
	\centering
	\includegraphics[width=0.3\textwidth]{\logoKeter}\bigpar
	{\LARGE \textbf{Proyecto Final}\par\vspace{0.2cm}}
	{\Huge\bfseries\textcolor{colorTitulo}{\nombrePractica}}
\end{titlepage}

% Indice
\tableofcontents

% Documentacion
\newpage
\textcolor{colorTitulo}{\section{Introducción Teórica}}
\textbf{Keter Vulnerability} es una página web diseñada para representar algunos
errores de configuración a la hora de crear aplicaciones web. El objetivo es disponer a
los usuarios de diferentes \textbf{retos} donde deberán usar sus conocimientos y su capacidad de
búsqueda a través de Internet para explotar dichas vulnerabilidades.\bigpar

\begin{figure}[hbt]
	\centerline{\includegraphics[width=1\textwidth]{img/vuln.png}}
	\caption[Vulnerability imagen]{Fuente \href{https://whiteknightit.com/2019/09/18/xeon-and-other-intel-cpus-hit-by-netcat-security-vulnerability/}{whiteknightit.com}}
	\bigpar
\end{figure}

El objetivo es ser capaces de concienciar a jóvenes \textbf{desarrolladores} del peligro que puede
llevar realizar páginas web sin conciencia sobre los errores. Al final veremos como el usuario
es un factor bastante peligroso para los desarrolladores y como nunca podemos confiar plenamente
en sus intenciones.\par
Es por ello que debemos protegernos con las últimas tecnologías. Mantenernos \textbf{actualizados y
	activos} será la clave para evitar cualquier error.\bigpar

El objetivo de esta página es su escalabilidad. Una vez terminado el proyecto la página pasará a ser código abierto,
con el objetivo de que la comunidad pueda aportar sus propios retos y módulos desde GitHub. Es por ello que gran parte
de los contenidos serán en inglés.
\newpage

Haremos uso de las últimas tecnologías para la realización de este proyecto:
\begin{itemize}
	\item \textbf{LaTex}: haremos uso de esta herramienta de texto mediante código para la creación de esta misma documentación.
	\item \textbf{NodeJS}: utilizaremos este paquete de recurso para la creación de la página web.
	\item \textbf{Mongo Atlas}: nos prooverá una base de datos principal para retos de NoSQL Injection.
	\item \textbf{HerokuApps}: será la página que hosteará nuestra aplicación.
	\item \textbf{GitHub}: allí subiremos el código y utilizaremos la función de GitHub Pages para crear una pequeña página web que muestre un breve resumen.
	\item \textbf{Docker}: crearemos un proyecto adicional con un contenedor en Docker para poder montar y utilizar nuestra aplicación en local.
\end{itemize}
\begin{figure}[hbt]
	\raggedright{\includegraphics[width=0.16\textwidth]{img/latex.png}}
	\raggedright{\includegraphics[width=0.16\textwidth]{img/node.png}}
	\raggedright{\includegraphics[width=0.16\textwidth]{img/mongo.png}}
	\raggedright{\includegraphics[width=0.16\textwidth]{img/heroku.png}}
	\raggedright{\includegraphics[width=0.16\textwidth]{img/github.png}}
	\raggedright{\includegraphics[width=0.16\textwidth]{img/docker.png}}
\end{figure}
\clearpage


\textcolor{colorTitulo}{\section{Introducción Técnica}}

\clearpage


\textcolor{colorTitulo}{\section{Vulnerabilidades utilizadas}}


\clearpage


\textcolor{colorTitulo}{\section{API}}

\clearpage


\textcolor{colorTitulo}{\section{Código}}

\clearpage


\textcolor{colorTitulo}{\section{Explotaciones}}

\clearpage


\textcolor{colorTitulo}{\section{Soluciones}}

\clearpage


\textcolor{colorTitulo}{\section{Conclusión}}

\clearpage


\end{document}